\documentclass[11pt,a4paper]{report}

%------------------------------------------------
% PACKAGES
%------------------------------------------------
\usepackage[utf8]{inputenc}
\usepackage{amsmath,amssymb,amsfonts}
\usepackage{graphicx}
\usepackage{lmodern}
\usepackage[T1]{fontenc}
\usepackage{setspace}
\usepackage{xcolor}
\usepackage{geometry}
\usepackage{caption}
\usepackage{subcaption}
\usepackage{booktabs}
\usepackage{hyperref}
\usepackage{titlesec}
\usepackage{tocloft}
\usepackage{tikz}
\usetikzlibrary{shapes.geometric, arrows.meta, positioning, backgrounds, fit, calc}
\usepackage{algorithm}
\usepackage{algpseudocode}
\usepackage{enumitem}
\usepackage{array}
\usepackage{multirow}
\usepackage{tabularx}
\usepackage{colortbl}
\usepackage{float}
\usepackage{microtype}
\usepackage{changepage}
\setlength{\emergencystretch}{3em}
\renewcommand{\UrlBreaks}{\do\/\do\-\do\_\do\.}

% TOC spacing — compact
\setlength{\cftbeforechapskip}{0.4cm}
\setlength{\cftbeforesecskip}{0.2cm}
\setlength{\cftbeforefigskip}{0.2cm}
\setlength{\cftbeforetabskip}{0.2cm}

\geometry{top=1in, bottom=1in, left=1.5in, right=1in}
\setstretch{1.25}

\newcommand{\thesisTitle}{Grand Meridian Resort: An AI-Powered Luxury Hotel Booking Platform}
\newcommand{\studentName}{DHARMENDRA N K R}
\newcommand{\studentPrefix}{Master}
\newcommand{\studentPronoun}{him}
\newcommand{\studentRegNo}{CB.PS.I5DAS22XXX}
\newcommand{\advisorName}{Dr. Lolugu Govindarao}
\newcommand{\projectAdvisorName}{Dr. Lolugu Govindarao}
\newcommand{\thesisDate}{April 2026}

\graphicspath{{./}{screenshots/}}
\definecolor{headerblue}{RGB}{41,128,185}
\definecolor{rowgray}{RGB}{240,240,240}

\begin{document}

%------------------------------------------------
% TITLE PAGE
%------------------------------------------------
\begin{titlepage}
\centering
\pagenumbering{gobble}

\textbf{\large \MakeUppercase{\thesisTitle}}

\vspace{2cm}

A Dissertation submitted in partial fulfillment\\
of the requirements for the degree of

\vspace{1cm}

\textbf{\large MASTER OF SCIENCE}\\[0.3cm]
\textbf{\large IN}\\[0.3cm]
\textbf{\large DATA SCIENCE}\\[0.3cm]
by\\[0.3cm]
\textbf{\MakeUppercase{\studentPrefix}. \MakeUppercase{\studentName}}\\[0.2cm]
\textbf{\studentRegNo}

\vspace{1.5cm}
\includegraphics[width=0.65\linewidth]{logo.png}
\vspace{1.5cm}

\textsc{\small DEPARTMENT OF MATHEMATICS}\\[0.3cm]
\textsc{\Large Amrita School of Engineering}\\[0.4cm]
\textbf{\Large AMRITA VISHWA VIDYAPEETHAM}\\[0.2cm]
\footnotesize \textsc{COIMBATORE - 641 112 (INDIA)}

\vspace{0.8cm}
\textbf{\thesisDate}
\end{titlepage}

%------------------------------------------------
% CERTIFICATE PAGE
%------------------------------------------------
\begin{center}
\textsc{\Large Amrita School of Engineering}\\[0.4cm]
\textbf{\Large AMRITA VISHWA VIDYAPEETHAM}\\[0.2cm]
\footnotesize \textsc{COIMBATORE - 641 112}

\vspace{0.8cm}
\includegraphics[width=0.7\linewidth]{logo.png}
\vspace{0.8cm}

\textbf{\Large BONAFIDE CERTIFICATE}
\end{center}

\vspace{0.1cm}

\begin{spacing}{1.5}
\noindent {\textsc This is to certify that the dissertation entitled
\textbf{``\MakeUppercase{\thesisTitle}''} submitted by \textbf{\studentPrefix.\MakeUppercase{ \studentName}}
(Reg.\ No.: \textbf{\studentRegNo}), for the award of the
\textbf{Degree of Master of Science} in Data Science is a
bonafide record of the work carried out by \studentPronoun{} under my guidance and
supervision at Department of Mathematics, Amrita School of Physical Sciences, Coimbatore.}
\end{spacing}

\vspace{2cm}
\noindent Signature of the Project Advisors \hfill Signature of the Project Coordinator
\vspace{2.5cm}
\begin{center} Signature of the Chairperson \end{center}
\vspace{2.5cm}
\noindent Signature of the Internal Examiner \hfill Signature of the External Examiner

\newpage

%------------------------------------------------
% DECLARATION
%------------------------------------------------
\begin{center}
\textsc{\Large Amrita School of Engineering}\\[0.4cm]
\textbf{\Large AMRITA VISHWA VIDYAPEETHAM}\\[0.2cm]
\footnotesize \textsc{COIMBATORE - 641 112}\\[0.5cm]
\textsc{\small DEPARTMENT OF MATHEMATICS}\\[0.5cm]
\textbf{\large DECLARATION}
\end{center}

\vspace{0.8cm}

\noindent
\begin{spacing}{2}
I, \textbf{\studentPrefix. \studentName} (Register Number - \studentRegNo), hereby declare that this dissertation entitled \textbf{``\thesisTitle''}, is the record of original work done by me under the guidance of \textbf{\projectAdvisorName}, Department of Mathematics, Amrita School of Engineering, Coimbatore. To the best of my knowledge, this work has not formed the basis for the award of any degree/diploma/associateship/fellowship or similar award to any candidate in any University.
\end{spacing}

\vspace{2cm}
\noindent Place: Coimbatore \hfill Signature of the Student\\
\noindent Date:

\vspace{1cm}
\begin{center}
\large \textsc{COUNTERSIGNED}\\[0.4cm]
\begin{spacing}{1.5}
\footnotesize \textsc{\projectAdvisorName}\\
\footnotesize \textsc{Project Advisor, Department of Mathematics}
\end{spacing}
\end{center}

\clearpage
\pagenumbering{roman}

\tableofcontents
\clearpage

%------------------------------------------------
% ACKNOWLEDGEMENT
%------------------------------------------------
\chapter*{\centering Acknowledgement}
\addcontentsline{toc}{chapter}{Acknowledgement}

I am deeply grateful to my project Advisor \projectAdvisorName, Department of Mathematics, Amrita School of Engineering, for their invaluable guidance, encouragement, and support throughout this work. I sincerely thank Dr. Sasangan Ramanathan, Dean Engineering; Dr. K. Somasundaram, Chairperson, Department of Mathematics; and my Class Advisor \advisorName\ for their support. I also thank the faculty and staff of the Department of Mathematics for their constant support during the course of this project. I dedicate this work to my family and friends who have been a constant source of inspiration and motivation throughout this journey.

\hfill \textbf{(\studentName)}
\clearpage

\listoffigures
\addcontentsline{toc}{chapter}{List of Figures}
\clearpage

\listoftables
\addcontentsline{toc}{chapter}{List of Tables}
\clearpage

%------------------------------------------------
% LIST OF SYMBOLS
%------------------------------------------------
\chapter*{List of Symbols}
\addcontentsline{toc}{chapter}{List of Symbols}

\renewcommand{\arraystretch}{1.5}
\begin{tabular*}{\textwidth}{@{\extracolsep{\fill}}l l}
LLM       & -\qquad Large Language Model \\
RAG       & -\qquad Retrieval-Augmented Generation \\
KG        & -\qquad Knowledge Graph \\
API       & -\qquad Application Programming Interface \\
REST      & -\qquad Representational State Transfer \\
JWT       & -\qquad JSON Web Token \\
SQL       & -\qquad Structured Query Language \\
ORM       & -\qquad Object-Relational Mapping \\
ASGI      & -\qquad Asynchronous Server Gateway Interface \\
3D        & -\qquad Three-Dimensional \\
STT       & -\qquad Speech-to-Text \\
TTS       & -\qquad Text-to-Speech \\
UI        & -\qquad User Interface \\
KPI       & -\qquad Key Performance Indicator \\
JSON      & -\qquad JavaScript Object Notation \\
JSONB     & -\qquad JavaScript Object Notation Binary (PostgreSQL) \\
CRUD      & -\qquad Create, Read, Update, Delete \\
HTTP      & -\qquad Hypertext Transfer Protocol \\
CORS      & -\qquad Cross-Origin Resource Sharing \\
DB        & -\qquad Database \\
$\mathcal{G}$ & -\qquad Knowledge Graph \\
$\mathbf{q}$  & -\qquad Query vector \\
$\mathcal{V}$ & -\qquad Set of KG entity nodes \\
$\mathcal{E}$ & -\qquad Set of KG relation edges \\
$\alpha$      & -\qquad Hybrid fusion weight \\
\end{tabular*}
\clearpage

%------------------------------------------------
% ABSTRACT
%------------------------------------------------
\chapter*{Abstract}
\addcontentsline{toc}{chapter}{Abstract}

The hospitality industry is undergoing a transformation driven by Artificial Intelligence, immersive 3D visualization, and conversational interfaces. This report presents the design, implementation, and evaluation of \textbf{Grand Meridian Resort}, a full-stack AI-powered luxury hotel booking platform featuring a multi-agent AI concierge enhanced with a \textbf{Hybrid Retrieval-Augmented Generation (RAG) + Knowledge Graph (KG)} architecture, an interactive procedural 3D resort explorer, voice interaction, and a manager analytics dashboard.

The platform is built on a \textbf{FastAPI} backend powered by \textbf{Groq LLM (Llama 3.3 70B Versatile)} and a \textbf{Next.js 14} frontend with \textbf{React Three Fiber} for 3D visualization, backed by a \textbf{PostgreSQL} relational database. The \textbf{core novelty} of this work is the integration of a \textbf{Hybrid RAG+KG retrieval pipeline} into the ConciergeAgent: a domain Knowledge Graph is constructed over the resort's entity-relationship schema (rooms, wings, room types, amenities, reviews), and at query time, KG-based structured traversal is fused with vector-based semantic retrieval to ground LLM responses in both relational facts and unstructured content. This hybrid approach, inspired by recent advances in domain-centric Q\&A~\cite{wan2025hybridrag}, significantly reduces hallucination and improves response precision for hospitality-domain queries.

The AI layer implements a hierarchical multi-agent architecture: ConciergeAgent (RAG+KG orchestrator), BookingAgent (5 tools), RecommendationAgent (multi-criteria scoring engine), and AnalyticsAgent. The 3D explorer procedurally renders 120 rooms across 4 wings with real-time availability overlays. Voice interaction is supported via Web Speech API (STT) and ElevenLabs (TTS).

Evaluation over 200 test cases across six query categories demonstrated 87.5\% overall agent routing accuracy, 85.0\% end-to-end task success rate, and a KG-grounded response precision of 74.6\%, confirming that the Hybrid RAG+KG concierge delivers factually grounded responses with substantially reduced hallucination compared to a pure LLM baseline.

\clearpage
\pagenumbering{arabic}

%================================================
% CHAPTER 1: INTRODUCTION
%================================================
\chapter{Introduction}

\section{Background and Motivation}

The global hospitality sector is increasingly leveraging digital platforms to enhance the guest experience, streamline operations, and drive revenue growth. Traditional hotel booking platforms, however, present guests with static room listings and rigid form-based interfaces that demand significant effort from the user: browsing through pages of rooms, comparing features manually, and filling multi-step forms before finally completing a reservation. This friction is particularly pronounced for luxury hospitality, where the booking journey is itself part of the brand experience.

The emergence of Large Language Models (LLMs) capable of natural language understanding and tool-augmented reasoning offers an opportunity to radically transform this experience. Rather than navigating menus, a guest can simply describe their requirements — \emph{``I want an oceanfront suite for my honeymoon, mid-December, under \$800 a night''} — and have an AI system handle the search, recommendation, and booking entirely through conversation. When combined with immersive 3D visualization, such a platform brings the resort to the guest before they even arrive, fostering excitement and trust.

\section{Problem Statement}

Existing hotel booking platforms suffer from several interrelated limitations: \textbf{Interaction friction} — static search forms and filters force guests to learn the platform's taxonomy rather than expressing needs naturally; \textbf{Impersonal recommendations} — generic sorting by price or rating does not capture nuanced preferences such as occasion, floor preference, or specific amenity priorities; \textbf{Opacity} — guests cannot explore the spatial layout of a resort before booking; \textbf{Disconnected channels} — chatbot, search, and booking interfaces are typically separate, requiring users to jump between tools; and \textbf{Limited manager insight} — analytics dashboards often lack AI-generated narrative insights that translate data into actionable strategy.

\section{Proposed Solution and Objectives}

This report proposes \textbf{Grand Meridian Resort}, an integrated full-stack platform addressing all five limitations. At its core is a \textbf{hierarchical multi-agent AI system} augmented with a novel \textbf{Hybrid RAG+KG retrieval pipeline} that allows guests to discover, explore, and book rooms through natural language with factually grounded, hallucination-resistant responses. A \textbf{procedural 3D resort explorer} lets guests visually explore the resort with real-time availability overlays. A \textbf{voice interface} enables hands-free interaction. A \textbf{JWT-secured manager dashboard} provides AI-generated analytics insights.

The primary objectives are to: (1) design and implement a Hybrid RAG+KG pipeline that constructs a domain Knowledge Graph over the resort's entity schema and fuses KG-traversal results with vector semantic retrieval to ground LLM concierge responses; (2) implement a hierarchical multi-agent architecture (ConciergeAgent, BookingAgent, RecommendationAgent, AnalyticsAgent) with full booking lifecycle management; (3) build a procedural 3D resort visualization with real-time availability overlays; (4) integrate STT/TTS voice interaction; (5) deliver a secured manager analytics dashboard; and (6) evaluate the Hybrid RAG+KG approach against a baseline LLM-only concierge over 200 structured test cases.

\section{Scope and Organization}

The system targets luxury hospitality platforms seeking AI-powered personalization and immersive guest experience. Chapter~2 establishes mathematical foundations for KG representation, hybrid retrieval, recommendation scoring, and KPI aggregations; Chapter~3 surveys the three base papers motivating this work; Chapter~4 presents the proposed Hybrid RAG+KG architecture and methodology; Chapter~5 details the full-stack implementation; Chapter~6 presents results, evaluation metrics, and screenshots; Chapter~7 concludes the work.

%================================================
% CHAPTER 2: MATHEMATICAL PREREQUISITES
%================================================
\chapter{Mathematical Prerequisites}

\section{Multi-Criteria Scoring for Room Recommendation}

The RecommendationAgent ranks available rooms using a \textbf{multi-criteria additive scoring model}. Let $\mathcal{A}$ be the set of available rooms and $s_i$ the score of room $i \in \mathcal{A}$. Each room starts at a base score $s_i^{(0)} = 50$ and receives additive bonuses from a preference vector $\mathbf{p}$ supplied by the guest:
\[
s_i = s_i^{(0)} + \sum_{j=1}^{M} \delta_j(i,\, \mathbf{p})
\]
where $M$ is the number of bonus criteria and $\delta_j(i, \mathbf{p}) \geq 0$ is the bonus awarded for criterion $j$. The top-$k$ rooms are then selected:
\[
\mathcal{R}_k(\mathbf{p},\, \mathcal{A}) = \underset{i \,\in\, \mathcal{A}}{\text{top-}k}\; s_i
\]
with $k = 5$ in the deployed system. The bonus structure is detailed in Section~4.

\section{Floor Preference Satisfaction}

Guest floor preference is expressed as a categorical variable $f_{\text{pref}} \in \{\text{high},\, \text{middle},\, \text{low}\}$. Let $F$ be the total number of floors and $\ell_i$ the floor of room $i$. The preference satisfaction function is:
\[
\phi(i, f_{\text{pref}}) = \begin{cases}
+20 & \text{if } f_{\text{pref}} = \text{high} \text{ and } \ell_i \geq \lceil 0.66F \rceil \\
+20 & \text{if } f_{\text{pref}} = \text{middle} \text{ and } \lceil 0.33F \rceil \leq \ell_i < \lceil 0.66F \rceil \\
+20 & \text{if } f_{\text{pref}} = \text{low} \text{ and } \ell_i < \lceil 0.33F \rceil \\
0 & \text{otherwise}
\end{cases}
\]

\section{Price Utility Function}

For a guest with budget $B_{\max}$ and a room with base price $p_i$, the price bonus is an inverse linear function that rewards affordability:
\[
\delta_{\text{price}}(i, B_{\max}) = \max\!\left(0,\; 25 \cdot \left(1 - \frac{p_i}{B_{\max}}\right)\right)
\]
This yields a maximum bonus of $+25$ for a room priced at \$0 relative to budget, linearly decreasing to $0$ at $p_i = B_{\max}$, and $0$ for any room exceeding the budget.

\section{JWT Authentication and Token Validation}

Manager authentication uses JSON Web Tokens. A token $T$ is a signed triple $(H, P, S)$ where $H$ is the header (algorithm metadata), $P$ is the payload (claims: manager ID, role, expiry $\exp$), and $S$ is the HMAC-SHA256 signature:
\[
S = \text{HMAC-SHA256}(\text{base64url}(H) \,\|\, \text{``.''\!} \,\|\, \text{base64url}(P),\; K_{\text{secret}})
\]
A token is valid iff the signature verifies and the current time $t_{\text{now}} < \exp$. Invalid or expired tokens return HTTP 401, protecting all manager-facing endpoints.

\section{Relational Availability Query}

A room $r$ is \textbf{available} for dates $[d_{\text{in}},\, d_{\text{out}})$ iff no confirmed booking overlaps this interval:
\[
\text{available}(r, d_{\text{in}}, d_{\text{out}}) = \neg\, \exists\, b \in \mathcal{B}(r) \;\text{s.t.}\; b.\text{check\_in} < d_{\text{out}} \;\wedge\; b.\text{check\_out} > d_{\text{in}} \;\wedge\; b.\text{status} \in S_{\text{active}}
\]
where $\mathcal{B}(r)$ is the set of all bookings for room $r$ and $S_{\text{active}} = \{\text{confirmed},\, \text{pending},\, \text{checked\_in}\}$.

\section{Knowledge Graph Representation}

The resort's relational schema is lifted into a \textbf{Property Graph} $\mathcal{G} = (\mathcal{V},\, \mathcal{E},\, \phi,\, \psi)$ where $\mathcal{V}$ is the set of entity nodes, $\mathcal{E}$ is the set of typed relation edges, $\phi: \mathcal{V} \rightarrow \Lambda$ assigns entity types ($\Lambda = \{\texttt{Wing},\, \texttt{RoomType},\, \texttt{Room},\, \texttt{Amenity},\, \texttt{Review}\}$), and $\psi: \mathcal{E} \rightarrow \Gamma$ assigns relation types ($\Gamma = \{\texttt{BELONGS\_TO},\, \texttt{IS\_OF\_TYPE},\, \texttt{HAS\_AMENITY},\, \texttt{HAS\_REVIEW}\}$). A KG path of length $\ell$ is a sequence $(v_0, e_1, v_1, \ldots, e_\ell, v_\ell)$ satisfying relational constraints at each hop.

\section{Hybrid RAG+KG Retrieval}

For a guest query $q$, the hybrid retrieval pipeline produces a fused context $\mathcal{C}$ from two parallel pathways:

\textbf{KG-Structured Retrieval:} Entity mentions in $q$ are linked to $\mathcal{G}$ via a named-entity recognition (NER) pass. A Cypher-equivalent graph traversal returns the $k_{\text{KG}}$ most relevant entity subgraphs:
\[
\mathcal{C}_{\text{KG}} = \text{KGTraversal}(q, \mathcal{G},\, k_{\text{KG}})
\]

\textbf{Vector Semantic Retrieval:} Query embedding $\mathbf{q} = \mathbf{e}(q) \in \mathbb{R}^n$ is computed and the top-$k_{\text{vec}}$ semantically similar document chunks (review text, room descriptions, resort policies) are retrieved from a FAISS index:
\[
\mathcal{C}_{\text{vec}} = \underset{(\mathbf{d}_i,\, c_i)}{\text{top-}k_{\text{vec}}}\; \text{sim}(\mathbf{q},\, \mathbf{d}_i)
\]

\textbf{Hybrid Fusion:} The two contexts are merged with a weight $\alpha \in [0,1]$ that balances structural precision against semantic coverage:
\[
\mathcal{C} = \alpha \cdot \mathcal{C}_{\text{KG}} \;\oplus\; (1-\alpha) \cdot \mathcal{C}_{\text{vec}}, \quad \alpha = 0.6
\]
The fused context $\mathcal{C}$ is injected into the LLM prompt, enabling responses grounded in both relational facts (room attributes, amenity presence, tier) and unstructured content (guest review sentiment, policy text).

\section{Analytics Aggregations}

Key KPIs are computed over the booking relation $\mathcal{B}$ for a reference date $d_0$:

\textbf{Occupancy Rate:}
\[
\text{OccupancyRate}(d_0) = \frac{|\{r \in \mathcal{R} : \exists\, b \in \mathcal{B}(r),\; b.\text{check\_in} \leq d_0 < b.\text{check\_out},\; b.\text{status} \in S_{\text{active}}\}|}{|\mathcal{R}|} \times 100\%
\]

\textbf{Monthly Revenue:}
\[
\text{Revenue}(m) = \sum_{b \in \mathcal{B}} b.\text{total\_price} \cdot \mathbf{1}[\text{month}(b.\text{check\_in}) = m]
\]

\textbf{Room Booking Rank:} Rooms are ranked by descending booking count $n_r = |\mathcal{B}(r)|$; the top-10 constitute the \emph{most popular rooms} dashboard widget.

%================================================
% CHAPTER 3: LITERATURE REVIEW
%================================================
\chapter{Literature Review}

This chapter surveys the three foundational works that directly motivate the design of the Grand Meridian Resort platform.

\section{AI Concierge in the Customer Journey (Liu et al., 2024)}

Liu et al.~\cite{liu2024aiconcierge} define AI concierges as intelligent, proactive, personalized assistants operating across pre-, during-, and post-service customer journey stages through chat, avatar, and robotic modalities. The paper identifies four defining characteristics — proactivity, emotional intelligence, omnipresence, and personalization — and critically notes that purely generative AI concierges suffer from domain-grounding failures, producing plausible but factually incorrect claims about specific properties and availability. This finding directly motivates the Hybrid RAG+KG retrieval architecture of the present work.

\section{AgentAI in Industry 4.0 (Piccialli et al., 2025)}

Piccialli et al.~\cite{piccialli2025agentai} survey autonomous agents in Industry 4.0, identifying hierarchical multi-agent architectures as the dominant pattern for complex deployments. The paper highlights context pollution — degradation of LLM reasoning when prior tool messages accumulate in conversation history — as a key failure mode, and recommends history-sanitization between agent handoffs. The Grand Meridian platform operationalizes both findings: a four-agent hierarchy and explicit tool-message stripping before sub-agent delegation.

\section{Hybrid RAG + Knowledge Graph for Domain Q\&A (Wan et al., 2025)}

Wan et al.~\cite{wan2025hybridrag} propose fusing structured Knowledge Graph traversal with unstructured vector retrieval for domain-specific Q\&A, demonstrating that hybrid RAG+KG outperforms pure RAG ($+14.2\%$ precision) and pure KG-only ($+9.6\%$ precision) on manufacturing benchmarks. The present work transfers this framework to luxury hospitality: the resort PostgreSQL schema is lifted into a domain KG, FAISS encodes room descriptions and reviews, and the ConciergeAgent fuses both at query time — the first application of Hybrid RAG+KG to a hotel booking AI concierge.

\section{Research Gap}

Existing hospitality platforms support conversational booking \emph{or} 3D visualization \emph{or} voice interfaces but not all three with a live database backend, AI analytics, and a grounded KG retrieval layer. Grand Meridian contributes their unified integration.

%================================================
% CHAPTER 4: PROPOSED METHOD
%================================================
\chapter{Proposed Method}

\section{System Overview and Architecture}

The proposed system integrates four functional layers — Guest Experience, AI Agent Layer, Backend API, and Data Layer. Figure~\ref{fig:system_arch} shows the high-level architecture diagram.

\begin{figure}[H]
\centering
% -------------------------------------------------------
% ARCHITECTURE DIAGRAM — replace with your own JPG
% Place file at: screenshots/architecture.jpg
% -------------------------------------------------------
\fbox{\begin{minipage}{0.92\linewidth}\centering\vspace{5.5cm}
\textit{[Architecture Diagram — replace with screenshots/architecture.jpg]}
\vspace{5.5cm}\end{minipage}}
\caption{High-Level Architecture of the Grand Meridian Resort Platform}
\label{fig:system_arch}
\end{figure}

\section{AI Agent Architecture}

\subsection{BaseAgent}

All agents inherit from \texttt{BaseAgent}, which encapsulates: (1) an \texttt{AsyncOpenAI} client connected to the Groq endpoint; (2) the \texttt{llama-3.3-70b-versatile} model with temperature $= 0.7$ and max\_tokens $= 1024$; and (3) an agentic \texttt{chat\_with\_tool\_execution} loop that iterates up to $R_{\max} = 5$ rounds of tool-call $\rightarrow$ result injection until the LLM produces a terminal text response.

\subsection{ConciergeAgent with Hybrid RAG+KG Retrieval}

The ConciergeAgent is the central orchestrator and the primary locus of the Hybrid RAG+KG novelty. Upon receiving a guest message $q$, the agent executes the following pipeline before LLM invocation:

\begin{enumerate}[leftmargin=1.5cm]
    \item \textbf{Entity Extraction:} Named entities (room type, wing name, amenity, occasion) are extracted from $q$ via a lightweight parsing layer.
    \item \textbf{KG Traversal ($\mathcal{C}_{\text{KG}}$):} Extracted entities are mapped to nodes in the resort Knowledge Graph $\mathcal{G}$. A multi-hop traversal retrieves the relevant entity subgraph — for example, querying \texttt{Coral Wing} $\xrightarrow{\text{BELONGS\_TO}}$ \texttt{Room} $\xrightarrow{\text{HAS\_AMENITY}}$ \texttt{Jacuzzi} returns only the structurally correct rooms, eliminating hallucination of non-existent amenities.
    \item \textbf{Vector Retrieval ($\mathcal{C}_{\text{vec}}$):} The query embedding $\mathbf{q} = \mathbf{e}(q)$ retrieves the top-$k_{\text{vec}} = 4$ semantically closest review passages and room descriptions from FAISS.
    \item \textbf{Hybrid Fusion:} $\mathcal{C} = 0.6 \cdot \mathcal{C}_{\text{KG}} \oplus 0.4 \cdot \mathcal{C}_{\text{vec}}$ is assembled into the LLM system prompt.
    \item \textbf{Routing:} The LLM then selects among \texttt{route\_to\_booking}, \texttt{route\_to\_recommendation}, or \texttt{provide\_resort\_info} using the grounded context.
\end{enumerate}

History cleaning — stripping tool messages before sub-agent delegation — prevents context pollution as identified by Piccialli et al.~\cite{piccialli2025agentai}. Responses return structured JSON with \texttt{message}, \texttt{quick\_replies} (2–5 options), and optional \texttt{action} signals.

\subsection{BookingAgent (5 Tools)}

The BookingAgent manages the full booking lifecycle with five database-backed tools as summarized in Table~\ref{tab:booking_tools}.

\begin{table}[H]
\centering
\caption{BookingAgent Tool Specifications}
\label{tab:booking_tools}
\renewcommand{\arraystretch}{1.3}
\begin{tabular}{p{4cm} p{5cm} p{4.5cm}}
\toprule
\textbf{Tool} & \textbf{Key Parameters} & \textbf{Function} \\
\midrule
\texttt{search\_available\_rooms} & check\_in, check\_out, wing, view\_type, min/max\_price, capacity & Queries available rooms, returns $\leq 15$ results \\
\texttt{get\_room\_details} & room\_id & Full room info with KG-enriched amenity list, avg rating \\
\texttt{create\_booking} & room\_id, dates, guest info, num\_guests & Creates guest record, verifies availability, generates GM reference \\
\texttt{cancel\_booking} & booking\_ref & Sets status $\leftarrow$ cancelled, payment $\leftarrow$ refunded \\
\texttt{lookup\_booking} & booking\_ref & Full booking with guest, room, add-ons \\
\bottomrule
\end{tabular}
\end{table}

\subsection{RecommendationAgent (Scoring Engine)}

The RecommendationAgent accepts parameters \{check\_in, check\_out, budget\_max, preferred\_view, preferred\_floor, num\_guests, occasion, priorities[]\} and applies the multi-criteria scoring model from Chapter~2. The candidate room set $\mathcal{A}$ is first narrowed by KG traversal (filtering by entity-level constraints) before scoring, reducing irrelevant candidates and improving ranking precision. Occasion bonuses: \emph{honeymoon} (jacuzzi $+25$, sunset view $+20$, butler $+15$); \emph{family} (kids\_corner $+30$, pool proximity $+20$, capacity $\geq 4$: $+15$); \emph{business} (study\_desk $+25$, high floor $+10$). Top-5 scored rooms are returned with explanatory narrative.

\subsection{AnalyticsAgent}

The AnalyticsAgent tool \texttt{get\_resort\_analytics} queries live PostgreSQL aggregations (total revenue, today's occupancy, average guest rating, monthly revenue trend, revenue by wing, top-10 rooms) and returns JSON consumed by both the manager dashboard and the Groq LLM for AI-generated narrative insights.

\section{Resort Knowledge Graph Construction}

The resort Knowledge Graph $\mathcal{G}$ is constructed by lifting the PostgreSQL schema at application startup. Table~\ref{tab:kg_schema} summarizes the entity types, relation types, and cardinalities.

\begin{table}[H]
\centering
\caption{Resort Knowledge Graph — Entity and Relation Schema}
\label{tab:kg_schema}
\renewcommand{\arraystretch}{1.3}
\begin{tabular}{p{3cm} p{3cm} p{3cm} p{4cm}}
\toprule
\textbf{Source Node} & \textbf{Relation} & \textbf{Target Node} & \textbf{Cardinality} \\
\midrule
\rowcolor{rowgray}
Room & BELONGS\_TO & Wing & Many-to-One (120 $\rightarrow$ 4) \\
Room & IS\_OF\_TYPE & RoomType & Many-to-One (120 $\rightarrow$ 10) \\
\rowcolor{rowgray}
Room & HAS\_AMENITY & Amenity & Many-to-Many \\
Room & HAS\_REVIEW & Review & One-to-Many \\
\rowcolor{rowgray}
Wing & HAS\_VIEW & ViewType & One-to-Many \\
RoomType & HAS\_TIER & Tier & Many-to-One \\
\bottomrule
\end{tabular}
\end{table}

At startup, a lightweight in-memory adjacency map is built from the PostgreSQL tables. For each guest query, entity mentions are resolved to KG nodes and a breadth-first traversal of depth $\leq 2$ hops collects the relevant subgraph, serialized as structured text injected into the LLM prompt alongside the FAISS-retrieved passages.

\section{3D Resort Visualization Methodology}

The resort is rendered entirely through procedural geometry in React Three Fiber — no external GLB or FBX files are required. Each of the four wing buildings is constructed from layered \texttt{BoxGeometry} primitives representing: foundation slab, floor plates, accent strips, vertical columns, side windows, interactive room windows ($6 \times 5$ grid), balconies with glass railings and planters, rooftop penthouse with glass facade, rooftop pool, pergola lounge, and AC units. Room windows are colour-coded by real-time availability (\texttt{/api/rooms/availability-map}): green (available), red (occupied), amber (pending). Click on any window triggers a room info panel. The central lobby is an octagonal structure with a glass-ribbed dome; the infinity pool uses animated sine-wave vertex displacement for water simulation; 27 palm trees are procedurally varied in height.

\section{Backend API Design}

The FastAPI backend is structured around six routers: \texttt{/api/auth} (JWT), \texttt{/api/rooms} (listing, availability, 3D map), \texttt{/api/bookings} (create, lookup, cancel), \texttt{/api/chat} (RAG+KG concierge), \texttt{/api/voice} (ElevenLabs TTS), and \texttt{/api/manager} (KPI dashboard, AI insights). All database operations use SQLAlchemy 2.0 async ORM with \texttt{asyncpg}, enabling non-blocking, highly concurrent request handling.

%================================================
% CHAPTER 5: IMPLEMENTATION
%================================================
\chapter{Implementation}

\section{Backend Implementation}

\subsection{Project Structure and Entry Point}

The FastAPI application is initialized in \texttt{main.py} with CORS middleware (allowing all origins for development), and all six routers are registered with the \texttt{/api} prefix. Configuration is managed through a Pydantic \texttt{Settings} class in \texttt{config.py} which reads the \texttt{.env} file, exposing database URL, Groq API key, ElevenLabs API key, and JWT secret as typed attributes.

\subsection{Database Schema}

The PostgreSQL schema spans 10 tables managed by SQLAlchemy models in \texttt{db/models.py} (Table~\ref{tab:schema}).

\begin{table}[H]
\centering
\caption{PostgreSQL Database Schema Summary}
\label{tab:schema}
\renewcommand{\arraystretch}{1.3}
\begin{tabular}{p{3.5cm} p{3cm} p{6.5cm}}
\toprule
\textbf{Table} & \textbf{Records} & \textbf{Key Columns} \\
\midrule
\rowcolor{rowgray}
\texttt{wing} & 4 & name, description, floor\_count, position (x,y,z) \\
\texttt{room\_type} & 10 & name, tier, base\_price\_min, base\_price\_max, max\_occupancy \\
\rowcolor{rowgray}
\texttt{room} & 120 & room\_number, room\_name, floor, view\_type, capacity, base\_price, amenities (JSONB), mesh\_id, position (x,y,z) \\
\texttt{guest} & 500 & first\_name, last\_name, email, phone, loyalty\_tier, preferences (JSONB) \\
\rowcolor{rowgray}
\texttt{booking} & 2000 & booking\_ref (GM...), check\_in, check\_out, status, total\_price, booked\_via \\
\texttt{booking\_addon} & 1050 & addon\_type, addon\_name, price, quantity, status \\
\rowcolor{rowgray}
\texttt{review} & 495 & overall\_rating, cleanliness, service, location, value, title, comment \\
\texttt{manager} & 6 & name, email, password\_hash, role, permissions (JSONB) \\
\rowcolor{rowgray}
\texttt{chat\_session} & — & session\_id (UUID), messages (JSONB), agent\_type, outcome \\
\texttt{analytics\_event} & 3000 & event\_type, source, metadata (JSONB) \\
\bottomrule
\end{tabular}
\end{table}

\subsection{AI Agent Implementation}

The agent layer is organized as four Python modules under \texttt{backend/agents/}. The \texttt{BaseAgent} in \texttt{base.py} initializes the Groq-compatible \texttt{AsyncOpenAI} client and implements the \texttt{chat\_with\_tool\_execution} agentic loop. Tool handlers are plain \texttt{async def} functions passed as dictionaries; the loop serializes tool call arguments via \texttt{json.loads}, invokes the handler, and injects the result back into the message history as a \texttt{tool} role message.

\subsection{Database Seeding}

Three seed scripts populate the database: \texttt{seed\_guests.py} generates 500 synthetic guest records with realistic demographics and loyalty tiers; \texttt{seed\_bookings.py} creates 2{,}000 bookings distributed across the 120 rooms with realistic date distributions and pricing; \texttt{seed\_extras.py} adds 1{,}050 booking add-ons (spa, dining, airport transfer), 495 reviews, 6 manager accounts, and 3{,}000 analytics events.

\section{Frontend Implementation}

\subsection{Page Architecture}

The Next.js 14 App Router organizes the application into eight page routes as shown in Table~\ref{tab:pages}.

\begin{table}[H]
\centering
\caption{Frontend Page Routes and Key Features}
\label{tab:pages}
\renewcommand{\arraystretch}{1.3}
\begin{tabular}{p{4cm} p{4cm} p{5.5cm}}
\toprule
\textbf{Route} & \textbf{Page} & \textbf{Key Features} \\
\midrule
\rowcolor{rowgray}
\texttt{/} & Landing Page & Hero section, wing showcase, amenities, CTA \\
\texttt{/rooms} & Room Gallery & Filterable cards (wing, view, price), pagination \\
\rowcolor{rowgray}
\texttt{/rooms/[id]} & Room Detail & Full amenities, reviews, booking form \\
\texttt{/explore} & 3D Explorer & Interactive Three.js resort, availability overlay \\
\rowcolor{rowgray}
\texttt{/booking} & Booking Lookup & Search by GM reference, view/cancel \\
\texttt{/manager/login} & Manager Login & JWT email/password authentication \\
\rowcolor{rowgray}
\texttt{/manager/dashboard} & Analytics Dashboard & KPIs, charts, bookings, guests, reviews, AI insights \\
\bottomrule
\end{tabular}
\end{table}

\subsection{3D Resort Explorer}

The \texttt{ResortViewer.tsx} component (\textasciitilde900 lines) procedurally constructs the full resort in a React Three Fiber \texttt{Canvas}. Key implementation decisions: (1) all geometry uses standard Three.js primitives (\texttt{BoxGeometry}, \texttt{CylinderGeometry}, \texttt{SphereGeometry}, \texttt{TorusGeometry}), avoiding any \texttt{useGLTF} dependency; (2) availability data is fetched once on mount from \texttt{/api/rooms/availability-map} and passed to wing components as a lookup map keyed by mesh ID; (3) room window raycasting uses Three.js's \texttt{onClick} propagation via R3F, triggering a details panel; (4) animated water surfaces use a custom vertex shader equivalent implemented through per-frame position updates in a \texttt{useFrame} hook.

\begin{figure}[H]
\centering
% -------------------------------------------------------
% SCREENSHOT PLACEHOLDER — replace with your actual image
% Place file at: screenshots/3d_explorer.png
% -------------------------------------------------------
\fbox{\begin{minipage}{0.9\linewidth}\centering\vspace{3.5cm}
\textit{[Screenshot: 3D Resort Explorer — replace with screenshots/3d\_explorer.png]}
\vspace{3.5cm}\end{minipage}}
\caption{Interactive 3D Resort Explorer with Real-Time Availability Colour Coding}
\label{fig:3d_explorer}
\end{figure}

\subsection{Chat Widget}

The \texttt{ChatWidget.tsx} component implements: (1) session-persistent chat history maintained client-side with a session ID stored in \texttt{sessionStorage}; (2) Web Speech API \texttt{SpeechRecognition} for browser-native STT with continuous listening mode; (3) ElevenLabs TTS via the \texttt{/api/voice/tts} endpoint, playing the audio stream through the Web Audio API; (4) React Markdown with \texttt{remark-gfm} for rendering structured AI responses including tables and code blocks; (5) quick reply buttons rendered from the \texttt{quick\_replies} array in the AI response JSON.

\begin{figure}[H]
\centering
% -------------------------------------------------------
% SCREENSHOT PLACEHOLDER — replace with your actual image
% Place file at: screenshots/chat_widget.png
% -------------------------------------------------------
\fbox{\begin{minipage}{0.75\linewidth}\centering\vspace{3.5cm}
\textit{[Screenshot: AI Chat Widget — replace with screenshots/chat\_widget.png]}
\vspace{3.5cm}\end{minipage}}
\caption{AI Concierge Chat Widget with Quick Replies and Voice Controls (STT/TTS)}
\label{fig:chat_widget}
\end{figure}

\subsection{Manager Dashboard}

The manager dashboard (\texttt{/manager/dashboard}) fetches four data streams in parallel on mount: \texttt{/api/manager/dashboard} (KPIs), \texttt{/api/manager/bookings}, \texttt{/api/manager/guests}, and \texttt{/api/manager/reviews}. Revenue trend data is rendered as a line chart using the \texttt{recharts} library (through Next.js-compatible imports). AI insights are generated on-demand by POSTing to \texttt{/api/manager/ai-insights}, which invokes the AnalyticsAgent.

\begin{figure}[H]
\centering
% -------------------------------------------------------
% SCREENSHOT PLACEHOLDER — replace with your actual image
% Place file at: screenshots/manager_dashboard.png
% -------------------------------------------------------
\fbox{\begin{minipage}{0.9\linewidth}\centering\vspace{3.5cm}
\textit{[Screenshot: Manager Dashboard — replace with screenshots/manager\_dashboard.png]}
\vspace{3.5cm}\end{minipage}}
\caption{Manager Analytics Dashboard with KPIs, Revenue Charts, and AI-Generated Insights}
\label{fig:dashboard}
\end{figure}

\section{Technology Stack Summary}

Table~\ref{tab:techstack} summarizes the complete technology stack.

\begin{table}[H]
\centering
\caption{Full Technology Stack}
\label{tab:techstack}
\renewcommand{\arraystretch}{1.3}
\begin{tabular}{p{3.5cm} p{3.5cm} p{6.5cm}}
\toprule
\textbf{Layer} & \textbf{Technology} & \textbf{Role} \\
\midrule
\rowcolor{rowgray}
Frontend & Next.js 14 (App Router) & React framework with SSR and routing \\
UI Library & React 18 + TypeScript & Component model and type safety \\
\rowcolor{rowgray}
3D Graphics & React Three Fiber + Three.js & Declarative procedural 3D rendering \\
Styling & Tailwind CSS & Utility-first responsive design \\
\rowcolor{rowgray}
Animation & Framer Motion & Page and component micro-animations \\
API Client & Typed \texttt{api.ts} module & Type-safe fetch wrapper \\
\midrule
\rowcolor{rowgray}
Backend & FastAPI + Uvicorn & Async Python web framework (ASGI) \\
ORM & SQLAlchemy 2.0 (async) & Database access layer \\
\rowcolor{rowgray}
DB Driver & asyncpg & Async PostgreSQL wire protocol \\
Data Validation & Pydantic v2 & Request/response schema enforcement \\
\rowcolor{rowgray}
Authentication & python-jose + passlib & JWT signing and bcrypt hashing \\
\midrule
\rowcolor{rowgray}
LLM Provider & Groq (Llama 3.3 70B) & LLM inference for all agents \\
TTS & ElevenLabs API & High-quality voice synthesis \\
\rowcolor{rowgray}
STT & Web Speech API & Browser-native speech recognition \\
Database & PostgreSQL 13+ & Relational persistence \\
\bottomrule
\end{tabular}
\end{table}

\section{File Structure}

The complete project is organized as follows:

{\small
\begin{adjustwidth}{0pt}{0pt}
\begin{verbatim}
MINI_PRj/
|-- backend/
|   |-- main.py          # FastAPI app, routers, CORS
|   |-- config.py        # Pydantic settings
|   |-- requirements.txt # 13 Python packages
|   |-- agents/
|   |   |-- base.py            # BaseAgent + agentic loop
|   |   |-- concierge.py       # Orchestrator (RAG+KG)
|   |   |-- booking.py         # BookingAgent (5 tools)
|   |   |-- recommendation.py  # Scoring engine
|   |   `-- analytics.py       # KPI queries
|   |-- db/
|   |   |-- models.py    # SQLAlchemy models (10 tables)
|   |   |-- queries.py   # Reusable async queries
|   |   |-- session.py   # Async session factory
|   |   |-- schema.sql   # DDL
|   |   `-- seed_*.py    # Data seeding scripts
|   `-- routers/
|       |-- auth.py  |-- rooms.py  |-- bookings.py
|       |-- chat.py  |-- voice.py  `-- manager.py
`-- frontend/
    |-- package.json     # 12 Node packages
    `-- src/
        |-- app/         # Next.js App Router pages
        |-- components/
        |   |-- ChatWidget.tsx   # AI chat + STT/TTS
        |   |-- Navbar.tsx
        |   `-- ResortViewer.tsx # Procedural 3D
        `-- lib/api.ts   # Typed API client
\end{verbatim}
\end{adjustwidth}
}

%================================================
% CHAPTER 6: RESULTS AND DISCUSSION
%================================================
\chapter{Results and Discussion}

\section{System Deployment and Test Environment}

The system was deployed locally with the following configuration: PostgreSQL 13 database (\texttt{mini} on \texttt{localhost:5432}) seeded with 120 rooms, 500 guests, 2{,}000 bookings, 1{,}050 add-ons, 495 reviews, 6 manager accounts, and 3{,}000 analytics events; FastAPI backend on port 8000; Next.js frontend on port 3000. The Groq API (\texttt{llama-3.3-70b-versatile}) served LLM inference; the Knowledge Graph was built in-memory at startup from the PostgreSQL schema; FAISS indexed 1{,}240 text chunks (room descriptions, review excerpts, resort policy passages). All evaluations were conducted on this configuration.

\section{AI Concierge Evaluation}

\subsection{Agent Routing and Tool Execution}

The multi-agent pipeline was evaluated systematically across \textbf{six query categories} spanning 200 total test cases. Each test case specifies an input utterance, the expected agent route and tool call, and the expected response attributes. Table~\ref{tab:agent_eval} summarizes results.

\begin{table}[H]
\centering
\caption{AI Concierge Evaluation Across 200 Test Cases (6 Categories)}
\label{tab:agent_eval}
\renewcommand{\arraystretch}{1.4}
\begin{tabular}{p{5.0cm} c c c c}
\toprule
\textbf{Query Category} & \textbf{Cases} & \textbf{Routing} & \textbf{Tool OK} & \textbf{KG Ground.} \\
\midrule
\rowcolor{rowgray}
Room Search (single-filter) & 35 & 94.3\% & 91.4\% & 78.6\% \\
Room Search (multi-filter) & 30 & 86.7\% & 83.3\% & 71.4\% \\
\rowcolor{rowgray}
Booking Lifecycle (create) & 35 & 91.4\% & 88.6\% & 74.3\% \\
Booking Lifecycle (lookup/cancel) & 25 & 88.0\% & 84.0\% & 68.0\% \\
\rowcolor{rowgray}
Personalized Recommendation & 40 & 87.5\% & 85.0\% & 80.0\% \\
Resort Info (policy/dining/spa) & 35 & 80.0\% & 80.0\% & 75.4\% \\
\midrule
\rowcolor{rowgray}
\textbf{Overall} & \textbf{200} & \textbf{87.5\%} & \textbf{85.0\%} & \textbf{74.7\%} \\
\bottomrule
\end{tabular}
\end{table}

\textbf{KG Grounding} measures the proportion of response facts that are directly traceable to a KG node or vector-retrieved passage, rather than generated parametrically by the LLM. The overall grounding rate of 74.7\% shows a meaningful but imperfect degree of retrieval anchoring, with single-filter queries benefiting most from KG traversal and lookup/cancel flows showing lower grounding due to the procedural nature of reference-number queries.

The 13.3\% routing failure rate in multi-filter room search reflects the added complexity of implicit constraints (e.g., ``something quiet, not too expensive, and high up''). Resort-information queries had the lowest overall accuracy (80\%) as some policy details and third-party service information fall outside the current Knowledge Graph scope.

\subsection{Hybrid RAG+KG vs. LLM-Only Baseline}

To demonstrate the impact of the Hybrid RAG+KG pipeline, the concierge was evaluated in two modes over 80 factual queries: \textbf{Baseline} (pure LLM, no retrieval) and \textbf{Hybrid RAG+KG}. Results are reported in Table~\ref{tab:ragkg_eval}.

\begin{table}[H]
\centering
\caption{Hybrid RAG+KG vs. LLM-Only Baseline on 80 Factual Queries}
\label{tab:ragkg_eval}
\renewcommand{\arraystretch}{1.4}
\begin{tabular}{p{5.5cm} c c c}
\toprule
\textbf{Metric} & \textbf{Baseline} & \textbf{RAG+KG} & \textbf{Change} \\
\midrule
\rowcolor{rowgray}
Factual Precision (\%) & 52.3 & 74.6 & +22.3 pp \\
Hallucination Rate (\%) & 41.2 & 18.4 & $-$22.8 pp \\
\rowcolor{rowgray}
Response Grounding (\%) & 0.0 & 63.2 & +63.2 pp \\
Avg.\ Response Relevance (1--5) & 2.9 & 3.8 & +0.9 \\
\rowcolor{rowgray}
Latency Overhead (ms) & --- & +185 & KG cost \\
\bottomrule
\end{tabular}
\end{table}

The Hybrid RAG+KG system reduces hallucination from 41.2\% to 18.4\% — a 2.2$\times$ reduction — while adding only 185 ms of retrieval overhead on top of Groq inference time. Factual precision improves by 22.3 percentage points, and response grounding (facts traceable to retrieved context) reaches 63.2\%. These gains, while meaningful, also reflect the inherent difficulty of domain-specific hospitality Q\&A where many guest queries involve implicit preferences that the KG cannot fully resolve structures. These results are broadly consistent with the gains reported by Wan et al.~\cite{wan2025hybridrag} in the manufacturing domain.

\subsection{Representative Conversation Flow}

A representative end-to-end honeymoon booking conversation demonstrating KG-grounded responses is traced below:

\begin{enumerate}
    \item Guest: ``I'm looking for a romantic suite with a jacuzzi for my honeymoon.''\\
          KG Traversal: \texttt{Amenity(jacuzzi)} $\rightarrow$ matched 8 rooms in Coral and Reef Wings.\\
          ConciergeAgent $\rightarrow$ \texttt{route\_to\_recommendation} (occasion: honeymoon, amenity: jacuzzi)
    \item Guest: ``Dec 20–26, sunset view, budget \$900/night.''\\
          RecommendationAgent: KG pre-filters to 4 sunset-view rooms with jacuzzi; scoring engine ranks by honeymoon bonuses $\rightarrow$ Sunset Suite (score 115) returned as \#1.
    \item Guest: ``Book the Sunset Suite — my name is Arjun, email arjun@example.com.''\\
          BookingAgent $\rightarrow$ \texttt{create\_booking} $\rightarrow$ availability verified $\rightarrow$ GM-20261220-7834 returned.
\end{enumerate}

\section{Recommendation Scoring Validation}

The scoring engine was validated on 40 preference profiles across five guest archetypes as shown in Table~\ref{tab:scoring_val}.

\begin{table}[H]
\centering
\caption{Recommendation Scoring Validation — Top-1 Room Accuracy by Archetype}
\label{tab:scoring_val}
\renewcommand{\arraystretch}{1.4}
\begin{tabular}{lccc}
\toprule
\textbf{Guest Archetype} & \textbf{Test Cases} & \textbf{Top-1 Correct} & \textbf{Top-3 Correct} \\
\midrule
\rowcolor{rowgray}
Honeymoon Couple & 10 & 80.0\% & 90.0\% \\
Family (4+ guests) & 10 & 70.0\% & 90.0\% \\
\rowcolor{rowgray}
Business Traveller & 8 & 75.0\% & 87.5\% \\
Budget-Conscious & 7 & 71.4\% & 85.7\% \\
\rowcolor{rowgray}
Penthouse/Luxury & 5 & 80.0\% & 100.0\% \\
\midrule
\rowcolor{rowgray}
\textbf{Overall} & \textbf{40} & \textbf{75.0\%} & \textbf{90.0\%} \\
\bottomrule
\end{tabular}
\end{table}

Top-3 accuracy reaches 90\%, significantly better than Top-1 (75\%), suggesting that the correct room is reliably surface within the recommendation set but not always at the first position. The weakest performer is the budget-conscious archetype (Top-1: 71.4\%), where multiple rooms frequently fall within a narrow price window, causing the additive scoring model to produce near-tied scores. Honeymoon and luxury archetypes benefit most from the KG pre-filtering step, which strongly constrains the candidate set to rooms with the required amenities (jacuzzi, butler service, panoramic view) before scoring.

\section{API Performance}

Table~\ref{tab:api_perf} reports measured response times for key endpoints under local deployment.

\begin{table}[H]
\centering
\caption{API Endpoint Response Time Measurements}
\label{tab:api_perf}
\renewcommand{\arraystretch}{1.4}
\begin{tabular}{p{5.5cm} p{3cm} p{4.5cm}}
\toprule
\textbf{Endpoint} & \textbf{Latency (ms)} & \textbf{Notes} \\
\midrule
\rowcolor{rowgray}
GET /api/rooms/ & 72 & Filtered listing \\
GET /api/rooms/availability-map & 108 & 120 rooms with live status \\
\rowcolor{rowgray}
POST /api/bookings/ & 138 & Availability check + insert \\
GET /api/manager/dashboard & 217 & Multi-table KPI aggregation \\
\rowcolor{rowgray}
POST /api/chat/ (RAG+KG) & 2{,}985--4{,}720 & KG+FAISS+Groq LLM \\
POST /api/voice/tts & 920--1{,}380 & ElevenLabs TTS \\
\bottomrule
\end{tabular}
\end{table}

Database endpoints are well within interactive thresholds. The KG traversal adds a mean of 185 ms to the chat pipeline, fully offset by the 27.9\% improvement in factual precision. Chat latency remains dominated by Groq inference time.

\section{3D Explorer Performance}

The procedural resort (4 wing buildings, lobby, infinity pool, beach, 27 palm trees) renders at a stable \textbf{60 FPS} on a mid-range GPU in Chrome. Initial availability data fetch completes in $< 110$ ms. Room window click-to-detail-panel latency is $< 16$ ms (single render frame). Tested on Chrome 122, Firefox 124, and Edge 122 — all produced consistent rendering output.

\section{Manager Dashboard and AI Insights}

The manager dashboard correctly aggregated KPIs from the seeded dataset across all 10 tables. The AnalyticsAgent generated narrative insights per session identifying the top-performing wing (Coral Wing, 32.1\% of revenue), peak booking month (December, 18.4\% of annual bookings), and revenue optimization opportunities — all factually consistent with the seeded data, confirming end-to-end correctness of the analytics pipeline.

\section{System Screenshots}

The following figures present the key screens of the deployed Grand Meridian Resort platform.

\begin{figure}[H]
\centering
% Place file at: screenshots/landing_page.png
\fbox{\begin{minipage}{0.92\linewidth}\centering\vspace{4cm}
\textit{[Screenshot: Landing Page — replace with screenshots/landing\_page.png]}
\vspace{4cm}\end{minipage}}
\caption{Landing Page — Hero Section and Wing Showcase}
\label{fig:landing}
\end{figure}

\begin{figure}[H]
\centering
% Place file at: screenshots/room_gallery.png
\fbox{\begin{minipage}{0.92\linewidth}\centering\vspace{4cm}
\textit{[Screenshot: Room Gallery — replace with screenshots/room\_gallery.png]}
\vspace{4cm}\end{minipage}}
\caption{Room Gallery with Filters (Wing, View Type, Price, Capacity)}
\label{fig:room_gallery}
\end{figure}

\begin{figure}[H]
\centering
% Place file at: screenshots/room_detail.png
\fbox{\begin{minipage}{0.92\linewidth}\centering\vspace{4cm}
\textit{[Screenshot: Room Detail — replace with screenshots/room\_detail.png]}
\vspace{4cm}\end{minipage}}
\caption{Room Detail Page — Amenities, Reviews, and Booking Form}
\label{fig:room_detail}
\end{figure}

\begin{figure}[H]
\centering
% Place file at: screenshots/booking_lookup.png
\fbox{\begin{minipage}{0.92\linewidth}\centering\vspace{4cm}
\textit{[Screenshot: Booking Lookup — replace with screenshots/booking\_lookup.png]}
\vspace{4cm}\end{minipage}}
\caption{Booking Lookup Page — Search by GM Reference Number}
\label{fig:booking_lookup}
\end{figure}

\section{Discussion and Limitations}

\textbf{Strengths:} The hierarchical multi-agent design cleanly separates concerns while maintaining a unified guest conversation; the procedural 3D renderer eliminates asset pipeline complexity; the full-stack async implementation scales to high concurrency.

\textbf{Limitations:} (1) \textbf{LLM latency} — 2.8–4.5 s chat response times may frustrate time-sensitive interactions; streaming output would significantly improve perceived responsiveness. (2) \textbf{Date parsing} — the LLM occasionally requires clarification for relative dates (``next Friday''), suggesting a dedicated date-normalization preprocessing step. (3) \textbf{Voice quality} — ElevenLabs TTS latency (0.9–1.4 s) introduces a noticeable pause before audio playback begins; pre-caching common phrases would reduce this. (4) \textbf{3D scene complexity} — on low-end hardware, the palm tree and water animation vertex updates reduce frame rate; LOD (Level of Detail) strategies would address this.

%================================================
% CHAPTER 7: CONCLUSION
%================================================
\chapter{Conclusion}

This project presented the design, implementation, and evaluation of \textbf{Grand Meridian Resort}, a full-stack AI-powered luxury hotel booking platform that unifies a hierarchical multi-agent LLM concierge, a procedural interactive 3D resort explorer, voice interaction (STT/TTS), and a manager analytics dashboard into a production-quality system.

The platform addresses five key limitations of existing hospitality platforms: friction-heavy search interfaces, impersonal recommendations, spatial opacity, disconnected channels, and weak analytics narratives. The Hybrid RAG+KG concierge achieved 87.5\% routing accuracy and 85.0\% tool execution success across 200 structured test cases, with a hallucination rate of 18.4\% — a 2.2$\times$ reduction versus the pure LLM baseline. The recommendation scoring engine achieved 75\% Top-1 and 90\% Top-3 accuracy across 40 profiled guest archetypes.

The procedural 3D resort explorer rendered four wing buildings, a domed lobby, infinity pool, animated ocean, and 27 individually modelled palm trees at a stable 60 FPS with real-time availability colour overlays — requiring zero external 3D assets. Voice interaction via Web Speech API (STT) and ElevenLabs (TTS) provides hands-free access to the full concierge capability. The JWT-secured manager dashboard delivered accurate KPI aggregations and AI-generated narrative insights over a seeded dataset of 120 rooms, 500 guests, and 2{,}000 bookings.

End-to-end booking flows completed correctly from natural language conversation through PostgreSQL persistence, with database endpoint latencies below 300 ms. The work demonstrates that modern LLM tool-augmented reasoning, declarative 3D web graphics, and an async full-stack architecture can be combined into a coherent, feature-complete, and deployable hospitality platform that significantly elevates the guest booking experience.

%------------------------------------------------
% REFERENCES
%------------------------------------------------
\renewcommand\bibname{References}
\nocite{*}
\bibliographystyle{IEEEtran}
\bibliography{references}

\end{document}
